\documentclass{article}
\usepackage{amsmath,amssymb,amsthm}
\newtheorem{theorem}{Theorem}
\newtheorem{lemma}{Lemma}
\begin{document}
\title{All direction multisets for the critical nine-degree-four profile}
\author{Near-extremal campaign, 2026-11-15}
\maketitle
\section{N01: Scope and actual-face input}
Let $G\subseteq Q_8$ be C4-free with683 edges and suppose one specified
global parity has degree profile $(0,0,0,0,9,71,45,2,1)$.
Let $R_1,\ldots,R_9\subseteq[8]$ be the entire retained direction sets
of its nine degree-four roots; repetitions are allowed.
The actual-face equality gate in q8\_equality\_profiles/main.tex,
with its declared local imports, gives $L_F\in\{0,2\}$ for all112
actual $Q_6$ faces. Every other specified root has degree at least five
and is never low on a $Q_6$ face. For each omitted coordinate pair
$\{a,b\}$, summing its four fixed-bit faces therefore gives
\[
       \#\{j:\{a,b\}\subseteq R_j\}\equiv0\pmod2.       \tag{1}
\]
\begin{theorem}
After one simultaneous coordinate permutation and a relabeling of the
nine roots, identify the eight coordinates with $\mathbb F_2^3$.
The nine direction sets are precisely
\[
 \{\{x:\ell(x)=1\}:0\ne\ell\in(\mathbb F_2^3)^*\}
                 \ \mathbin{\uplus}\ \{T,T\},          \tag{2}
\]
where $T$ is some four-element subset of the eight coordinates.
There are four possible types of $T$ under the group preserving the
seven displayed masks: $0\notin T$ with its complement in the seven
nonzero points a line or an independent triple, and $0\in T$ with
$T\setminus\{0\}$ a line or an independent triple.
\end{theorem}
This is a necessary classification of directions. It does not classify
the nine actual vertex positions, supply a683-edge completion, exclude
the whole degree profile, or prove the original upper682.

\section{N02: An even orthogonal code from the same masks}
For each coordinate$a$, form $u_a\in\mathbb F_2^9$ whose$j$th entry
records membership $a\in R_j$. Equation(1) says
$u_a\cdot u_b=0$ for distinct$a,b$. Summing (1) over $b\ne a$
counts $3\operatorname{wt}(u_a)$, so every $u_a$ is even and
$u_a\cdot u_a=0$ also. Their span$U$ is thus an even self-orthogonal
binary subspace. It has full support on all nine column coordinates,
since each $R_j$ has four elements.

The even subspace $E\subseteq\mathbb F_2^9$ has dimension eight.
Its dot-product radical is zero: the orthogonal complement of $E$
in the full space is the span of the all-one vector, which is odd.
The restricted form is alternating. If an isotropic subspace has
dimension$k<4$, its orthogonal complement in$E$ has dimension$8-k>k$;
choose a vector in that complement outside the current isotropic
subspace, and extend the latter by one dimension. Repeating gives an even self-orthogonal dimension-four
space $C\supseteq U$. Its support remains full.

\section{N03: Full-support dimension-four codes of length nine}
\begin{lemma}
Every full-support even self-orthogonal dimension-four code of length
nine is the direct sum of a length-two repetition code and a
length-seven dimension-three code all of whose nonzero words have
weight four. The weight-two word is unique.
\end{lemma}
\begin{proof}
A weight-eight word would omit one coordinate$r$. Orthogonality
with it and evenness of any codeword$z$ imply $z_r=0$, contrary to
full support. Suppose there were no weight-two word. The fifteen
nonzero words would have weights four or six. Every coordinate is
one on exactly eight codewords, so the sum of weights is72. Thus six
words have weight six and nine have weight four.

The zero sets of the six weight-six words are triples. Two distinct
such triples intersect in exactly one point: their corresponding
weight-six words intersect in four points, since their intersection
is even, at least three, and less than six. Any one triple meets five
others on only three points, so some three triples have a common
point. Their pairwise intersections are that same point; consequently
their binary sum has weight seven. The binary sum of the three
corresponding weight-six codewords is its complement, a weight-two
word. This contradiction proves the existence of a weight-two word$b$.

Orthogonality to$b$ makes its two entries equal in every codeword.
Subtracting$b$ when necessary decomposes $C$ as $\langle b\rangle$
plus an even self-orthogonal dimension-three code$D$ supported on the
other seven coordinates, with full support there. If$D$ contained a
weight-two word, stripping it in the same way would give a full-support
even self-orthogonal dimension-two code on five coordinates. Its three
nonzero words have total weight ten, hence weights two, four, four;
the two weight-four words intersect in three coordinates, contradicting
orthogonality. Thus$D$ has no weight-two word. A weight-six word of$D$
would force its omitted coordinate to be zero in every even codeword,
again contradicting full support. All seven nonzero words of$D$
therefore have weight four. This also proves uniqueness of$b$.
\end{proof}

\section{N04: Recovering the eight coordinate points}
Choose a basis for$D$, and label its seven columns by vectors
$p\in\mathbb F_2^3$ with multiplicities$m_p$. Each nonzero basis
combination has weight four, so its character sum is$7-2\cdot4=-1$;
the zero character sum is seven. Character orthogonality, obtained by
pairing $x$ with $x+e_i$ when the$i$th coefficient is nonzero, gives
$m_0=0$ and $m_p=1$ for every$p\ne0$. Thus the words of$D$ have
entries $a\cdot p$, $p\ne0$, for $a\in\mathbb F_2^3$.

Each of the eight actual rows$u_a\in C$ is consequently of the form
\[
              ((x\cdot p)_{p\ne0},\tau,\tau).
\]
Let$n_x$ count rows with first parameter$x$. Each of the first seven
original columns has weight four, because its original retained mask
has size four. Therefore $\sum_x n_x=8$ and
$\sum_x n_x(-1)^{x\cdot p}=0$ for every$p\ne0$. The same elementary
character inversion yields $n_x=1$ for all eight$x$. One coordinate
permutation identifies those rows with the eight points$x$.
The two remaining columns are identical, and their common support$T$
has size four. This proves(2). The change of column labels only
relabels roots; no root was moved or translated separately.

\section{N05: Four types and exact finite calibration}
The permutations preserving the seven nonzero linear-function masks
are exactly $\operatorname{GL}(3,2)$ on the seven nonzero coordinate
points, fixing zero. Indeed the induced permutation preserves the
three-dimensional mask space and therefore acts by an invertible
linear transformation; conversely every such transformation preserves
that space. A three-subset of the seven nonzero points is either a
line $\{a,b,a+b\}$ or a linearly independent triple. Each type is a
single orbit because a map between ordered independent pairs or bases
extends linearly. There are seven lines and28 independent triples.
According to whether zero is in$T$, use either its nonzero complement
or its three nonzero members. This gives the four types with sizes
7,28,7,28, respectively, accounting for all70 four-subsets.

The accompanying enumeration independently generates all even
self-orthogonal subspaces of dimensions zero through four without
quotienting by a group. The counts are $1,255,5355,11475,2295$; of
the last,1080 have full support. Every such space has the decomposition
proved above. This is a calibration of the elementary proof, not a
substitute for coverage of actual root positions or completions.
\end{document}
